% !PRINTED 509
\weberpage{509}
%<<CONT
unchanged by $Q$; and \gesperrt{secondly}, every function that remains
unchanged by $Q$ is, by \gesperrt{Lagrange}'s theorem, contained
in $\Omega'$. But these [by §. 149, a), b)] are the characteristic
marks of the \gesperrt{Galois} group in the field $\Omega'$.
In order, therefore, to reduce the \gesperrt{Galois} group from degree $p$ to degree $q$,
\gesperrt{one} root of an auxiliary equation of degree $j$
must be adjoined to $\Omega$.

Theorem 3. we can also express thus:

\gesperrt{The normal field $N = \Omega(\alpha, \alpha_1 \ldots \alpha_{m-1})$ is a field
of degree $p$ over $\Omega$ and of degree $q$ over $\Omega' = \Omega(\psi)$.}

To the lowering of the group by the adjunction of $\psi$ there corresponds,
as already follows from the general principles,
a splitting of the \gesperrt{Galois} resolvent. Let us assume
that (as in §. 145) $g(t) = 0$ is the \gesperrt{Galois} resolvent and
$\varrho$ one of its roots, and that by the permutations of $Q$,
$\varrho$ passes over into

\[
\varrho, \varrho_1, \varrho_2 \ldots \varrho_{q-1},
\]

and by the permutations of the subsidiary group $Q\pi_i$ into

% [sic] The source omits the comma after \varrho_{0,i} in this parallel list.
\[
\varrho_{0,i}\ \varrho_{1,i}, \varrho_{2,i} \ldots \varrho_{q-1,i},
\]

then each of the factors of $g(t)$

\[
g_i(t) = (t - \varrho_{0,i})(t - \varrho_{1,i}) \ldots (t - \varrho_{q-1,i})
\]

is unchanged by the substitutions of $Q$ and is therefore contained in $\Omega(\psi)$.

If we set

\[
g(t, \psi) = (t - \varrho)(t - \varrho_1) \ldots (t - \varrho_{q-1}),
\]

then it results, if we carry out in it a permutation from $Q\pi_i$,
that

\[
g_i(t) = g(t, \psi_i),
\]

and $g(t)$ is thus decomposable in the following manner:

\weq{6}{g(t) = g(t, \psi)\, g(t, \psi_1) \ldots g(t, \psi_{j-1}),}

in which each factor is of degree $q$, and no two of these
factors have a common divisor.

If $Q$ and $Q'$ are two different divisors of the group $P$, then
$Q$ and $Q'$ will have certain permutations in common, among
which the identical permutation is always found. It is
possible that this is the only common permutation of $Q$
and $Q'$, and in that case these two groups are called \gesperrt{coprime}.
But still more common elements may
%CONT>>
